\begin{figure}[htbp]
\centering
\begin{tikzpicture}[x=1cm,y=1cm]
\node[font=\small\sffamily\bfseries] at (2.5,1.1) {P1 的稳定业务结果};
\node[font=\small\sffamily\bfseries] at (10,1.1) {查询时的当前信息};
\node[booknode,text width=5.1cm] (fixed) at (2.5,-.7)
 {P1 成功，商品 paper-001\\金额 3000，订单 O1，扣款 D1\\D1 当次扣后余额 7000};
\node[booknode,text width=4.3cm] (now1) at (10,0) {第一次查询\\当前余额 7000};
\node[booknode,text width=4.3cm] (now2) at (10,-1.8) {另一笔购买以后再查询\\当前余额 4000};
\draw[bookarrow] (now1)--node[right,font=\scriptsize]{账户继续变化}(now2);
\node[booknode,text width=11cm] at (6.1,-3.4)
 {每次响应另有自己的 Date、request\_id 与追踪信息\\业务结果可保持稳定，查询元数据随本次通信产生};
\draw[bookarrow,dashed] (fixed.east)--node[above,font=\scriptsize]{一并返回}(now1.west);
\end{tikzpicture}
\caption{当次扣后余额属于原流水，当前余额属于查询时刻。区分字段的时间含义，才能解释同一购买的不同响应。}
\label{fig:ch02-result-time}
\end{figure}
